%%% -*- coding: utf-8 -*-
% !TeX root = main-slides.tex

\lecture{Debugging in Unity}{unity-debugging}
\section{Debugging in Unity}

\showtitlepage{Console, Breakpoints, Profiler}

%%% Every example is taken from SEE, so the sources named in the listings can
%%% be opened next to the slide.

\subsection{Console and Debugger}

\begin{frame}[fragile]{The Console, and How to Get More from It}
  \framelabel{slide:console}
  \begin{lstlisting}[basicstyle=\scriptsize\ttfamily]
Debug.Log("a plain message\n");
Debug.LogWarning("something is odd\n");
Debug.LogError("something is wrong\n");   // red, and with a stack trace

// Assets/SEE/Game/Operator/LabelOperator.cs -- the second argument
Debug.LogError("Could not retrieve portal for node label.\n", gameObject);

// Assets/SEE/Net/Util/Logger.cs -- a channel that can be switched off
public static bool InternalLoggingEnabled = true;
  \end{lstlisting}
  \vspace{1ex}

  \begin{itemize}
    \item That second argument of \texttt{Log*} is the part almost everybody forgets: it makes
      the entry clickable and selects that game object in the hierarchy. SEE
      passes it in a single place.
    \item A log line is not free. It builds a string and walks the stack, so
      it does not belong in code that runs every frame.
    \item An unexpected \texttt{LogError} fails a Unity test, even one that is
      otherwise green. If the error is the expected outcome, announce it with
      \texttt{LogAssert.Expect}.
  \end{itemize}
\end{frame}

\begin{frame}[fragile]{Fail Where the Cause Is}
  \begin{lstlisting}[basicstyle=\scriptsize\ttfamily]
// Assets/SEE/Extensions/Components.cs
public static T MustGetComponent<T>(this GameObject gameObject)
{
    if (!gameObject.TryGetComponent(out T component))
    {
        throw new InvalidOperationException(
            $"Couldn't find component '{typeof(T).GetNiceName()}' "
            + $"on game object '{gameObject.FullName()}'");
    }
    return component;
}
  \end{lstlisting}
  \vspace{1ex}

  A missing component that is quietly returned as null becomes a
  \texttt{NullReferenceException} elsewhere, a few frames later, in a stack
  that no longer mentions the cause. Naming the type and the object turns an
  afternoon of bisecting into reading one line.
  \vspace{1ex}

  The same idea, one level up: \texttt{AbstractNetAction} serializes itself to
  JSON inside \texttt{\#if UNITY\_EDITOR} before sending, so a net action
  holding an unserializable field is caught while you test it rather than in
  a session with others.
\end{frame}

\begin{frame}{Breakpoints, and What They Hide}
  \begin{itemize}
    \item Rider and Visual Studio attach to the running Editor, and SEE ships
      the packages for both. \texttt{Debug.Break()} pauses play mode from
      inside the code, which is the way to stop on a state you cannot express
      as a condition.
    \item A breakpoint stops the player loop as well. Anything timed --- a
      network timeout, an animation, a download --- is distorted the moment
      you stop, so what you observe afterwards is not what would have
      happened.
    \item Stepping through an asynchronous method walks the compiler's state
      machine. The call stack shows \texttt{MoveNext}, not the chain of
      callers you were thinking of, which is why the console is often faster
      than the debugger for a task that resumes on the wrong thread.
    \item A forgotten task leaves no trace in the console. \emph{Window
      $\rightarrow$ UniTask Tracker} lists every \texttt{UniTask} still
      pending, which is where an awaited-forever \texttt{WaitUntil} shows up.
  \end{itemize}
\end{frame}

\subsection{Profiling}

\begin{frame}{Measure, Then Change}
  \framelabel{slide:profiler}
  \begin{itemize}
    \item \emph{Window $\rightarrow$ Analysis $\rightarrow$ Profiler}, record
      while playing, then click the frame that is late. The two columns worth
      knowing first are CPU time and \emph{GC Alloc} --- the latter answers
      the allocation question raised in the lecture on Unity features.
    \item \emph{Deep Profile} records every managed call instead of the
      instrumented ones. It finds the method you did not suspect, and it
      distorts every timing while doing so, so use it to locate and not to
      measure.
    \item Editor numbers are not player numbers. To measure the real thing,
      build with \emph{Development Build} and \emph{Autoconnect Profiler} and
      profile the build.
    \item There is not a single profiler marker in SEE, and none is needed to
      start: Unity's own samples already carry the frame. A
      \texttt{ProfilerMarker} is what you add once you want your own row in
      that timeline.
  \end{itemize}
\end{frame}

\begin{frame}[fragile]{Two Time Scales}
  \begin{lstlisting}[basicstyle=\scriptsize\ttfamily]
// Assets/SEE/Utils/Performance.cs
Performance p = Performance.Begin("loading graph data");
...
p.End(true); // prints the action and the milliseconds it took

// Example usage in Assets/SEE/Game/CityRendering/GraphRenderer.cs
Performance.Begin($"Node layout {Settings.NodeLayoutSettings.Kind} ...");
... // Layout is calculated
p.End();
  \end{lstlisting}
  \vspace{1ex}

  \begin{itemize}
    \item The Profiler answers which frame was late. It cannot answer how long
      loading a city took, because that spans hundreds of frames and no single
      one of them is to blame.
    \item For those, SEE times whole operations and writes the result to the
      console, where it survives the session. A Profiler capture does not.
  \end{itemize}
  \vspace{1ex}

  \hint{\textit{Performance} is a class written by SEE developers.}
\end{frame}

\begin{frame}[fragile]{An Async Method Has No Row}
  \framelabel{slide:asyncprofiling}
  \begin{lstlisting}[basicstyle=\scriptsize\ttfamily]
// Assets/SEE/DataModel/DG/IO/GXL/GraphsReader.cs, in LoadAsync
foreach (string gxlPath in sortedGraphNames)
{
    await graphCreator.LoadAsync(...);
    await MetricImporter.LoadCsvAsync(graph, csvFilename);
}

// ... and what the Profiler shows of it: one entry per piece
PlayerLoop
  UniTaskLoopRunnerUpdate    // where UniTask pumps its continuations
    GraphsReader.<LoadAsync>d__N.MoveNext()   // the state machine
  \end{lstlisting}
  \vspace{0.5ex}

  \begin{itemize}
    \item The compiler cuts an \texttt{async} body at every \texttt{await} and
      makes a state machine of the pieces. Each piece is one call to
      \texttt{MoveNext} and they resume in different frames, so no row in the
      Hierarchy view adds up to what \texttt{LoadAsync} cost --- and the row's
      parent is whatever pumped the continuation, not the caller.
    \item Deep Profile is usually needed to see the pieces at all:
      compiler-generated methods are not among the ones Unity instruments.
  \end{itemize}
\end{frame}

\begin{frame}[fragile]{Elapsed Time Is Not CPU Time}
  \begin{lstlisting}[basicstyle=\scriptsize\ttfamily]
// Assets/SEE/Utils/Performance.cs -- a stopwatch survives every await
Performance p = Performance.Begin("loading graph data");
int frameBefore = Time.frameCount;
await reader.LoadAsync(directory, ...);
p.End(true);
Debug.Log($"spanned {Time.frameCount - frameBefore} frames\n");

marker.Begin();               // a profiler sample survives nothing
await SomethingSlow();        // wrong: End lands in a later frame
marker.End();
  \end{lstlisting}
  \vspace{0.5ex}

  \begin{itemize}
    \item The two are one number for a synchronous method and two unrelated
      ones for an async method: nine seconds of loading can be a fraction of a
      second of CPU, spread over hundreds of frames.
    \item Milliseconds and frames spanned together already route the question.
      One late frame is a case for the Profiler; hundreds of ordinary ones
      never are.
    \item A sample must begin and end on the same thread in the same frame, so
      instrument the pieces between the awaits, never the await itself.
  \end{itemize}
\end{frame}

\begin{frame}[fragile]{The Work You Cannot See}
  \framelabel{slide:offthread}
  \begin{lstlisting}[basicstyle=\scriptsize\ttfamily]
// Assets/SEE/Graphs/IO/GXL/GXLParser.cs
await UniTask.SwitchToThreadPool();
while (await Reader.ReadAsync())   // the whole parse, off the main thread
{
    ...
}

// what it takes to get a row in the timeline for that thread
Profiler.BeginThreadProfiling("SEE", "GXL parsing");
try { ... }
finally { Profiler.EndThreadProfiling(); }
  \end{lstlisting}
  \vspace{0.5ex}

  \begin{itemize}
    \item The Profiler samples the main thread, the render thread and Unity's
      job workers. A thread from the .NET thread pool is none of those, and
      that is where SEE does its heaviest loading: the GXL parser, the CSV
      metric importer, the report parsers, the Git providers.
    \item So \emph{I profiled the load and every frame looked idle} is the
      expected outcome rather than a mystery.
  \end{itemize}
\end{frame}

\begin{frame}{A Profiler That Sees Every Thread}
  \framelabel{slide:dottrace}
  \begin{itemize}
    \item dotTrace is JetBrains' profiler for .NET. It attaches to a process and
      samples every managed thread, so the thread-pool work of the previous
      slide shows up without a single marker being added to the code.
    \item It is free for students: the JetBrains educational license covers
      dotUltimate, and dotTrace is part of it. It is the same application you
      apply for to get Rider.
    \item Its Timeline type records what each thread did over time --- running,
      waiting for a lock, blocked in file I/O, collecting garbage --- which is
      the view an asynchronous load calls for.
    \item Division of labor: Unity's Profiler says which frame was late and what
      it allocated, dotTrace says where the managed time went across all
      threads. Neither replaces the other.
  \end{itemize}
\end{frame}

\begin{frame}[fragile]{Confining a Capture to One Operation}
  \begin{lstlisting}[basicstyle=\scriptsize\ttfamily]
// Assets/SEETests/Graphs/IO/TestGraphIO.cs
using (DeepProfiler.Capture("TestReadingRealBigGraph"))
{
    await LoadGraphAsync(path);
}

// Assets/SEE/Utils/DeepProfiler.cs -- when the two ends are apart
DeepProfiler.Start("loading city");
...
DeepProfiler.Save();
  \end{lstlisting}
  \vspace{0.5ex}

  \begin{itemize}
    \item Every method transitively called within the scope is profiled, none of
      them individually instrumented, and the snapshot covers that section
      alone instead of the whole session.
    \item The scope may enclose an \texttt{await}: it is disposed wherever the
      continuation resumes, hence it spans the frames the operation really took.
    \item With no profiler attached, these calls do nothing beyond one warning,
      which is why they may stay in the code.
  \end{itemize}
\end{frame}

\begin{frame}{Getting the Snapshot}
  \begin{enumerate}
    \item In dotTrace, under \emph{New Process Run}, add an entry whose Path is
      the Unity executable and whose Arguments carry
      \texttt{-projectpath} with this project. In the \emph{Advanced Options},
      set \emph{Control profiling} to \emph{Using API}.
    \item That one setting decides whether any of this works. With dotTrace's
      default, \emph{Manual control}, the API is inert: the console then reports
      that no profiler is attached and no snapshot is written.
    \item Run the entry, which starts Unity under the profiler, and trigger the
      code --- the test above from the Test Runner, say. The console reports the
      snapshot as saved and it appears in dotTrace's controller window.
  \end{enumerate}
  \vspace{1ex}

  \hint{Only one editor can run per project, so close a running one first. The
  profiling type hardly matters here; the Timeline type of the Unity preset
  works.}
\end{frame}

\begin{exerciseframe}{Which Tool?}
  \begin{enumerate}
    \item Loading a city takes nine seconds. Which measurement would you
      reach for, and why is the Profiler the wrong one?
    \item The frame rate drops, but only while a node is being dragged.
      Which tool now, and which column?
    \item You profile the loading of a code city with Deep Profile on and find
      no expensive method anywhere, although loading plainly takes seconds.
      Where did the time go?
  \end{enumerate}
\end{exerciseframe}

\begin{solutionframe}{Which Tool?}
  \begin{enumerate}
    \item A stopwatch around the operation, which is what
      \texttt{Performance.Begin} and \texttt{End} are for. Nine seconds is
      hundreds of frames, and the Profiler reports per frame, so it would show
      hundreds of unremarkable ones instead of the one number you want.
    \item The Profiler, recording while you drag, and the \emph{GC Alloc}
      column first: a drop that appears only during an interaction is usually
      something allocating per frame in the path that the interaction takes.
    \item Onto a thread the Profiler does not sample. The parsers begin with
      \texttt{await UniTask.SwitchToThreadPool()}, and a thread from the .NET
      thread pool gets no row until the code asks for one with
      \texttt{Profiler.BeginThreadProfiling} (\slideref{slide:offthread}).
  \end{enumerate}
\end{solutionframe}

%%% Local Variables:
%%% mode: latex
%%% TeX-master: "main-slides"
%%% mode: reftex
%%% mode: flyspell
%%% End:
